\section{边郡的账：军费、迁徙与地方军权}
\label{sec:eastern-han-frontier-fiscal}

东汉的西、北边郡不是雒阳危机之外的远景。羌地冲突反复延续，凉州、并州的军费、流民和安置不断累积；鲜卑的机动压力又使北方防务难以依靠一条固定防线。正史记录将领出征与“平定”比记录军粮账、逃户名册更连续，这个材料落差本身说明中枢必须为看不全、也难以直接支付的日常成本作决定。

\eventanchor{EH-0048-SOUTHERN-XIONGNU}{建武二十四年南匈奴日逐王比归附}，为边郡带来一个可结盟的政治伙伴，也带来安置部众、划定驻牧和供给首领的持续事务。归附年份可靠，部众数量、牧地边界与自主程度却不能由“内附”二字决定。到永初元年，\eventanchor{EH-0107-LIANG-QIANG-WAR}{凉州、并州一带的羌乱再起}，军费、迁徙与安置的压力又回到同一组郡县。部落称谓和各郡受害程度均有争议，但人群在战争中被迫改换土地、编户与保护者，是边防花费会留下社会后果的具体方式。

段颎在一六七至一六九年的进军能短期打通部分通道，不能证明冲突条件从此消失。斩获、迁徙和范围多来自胜方军政语言；“羌”“鲜卑”也不能被当成内部一致的政治意志。牧地、道路、粮运和地方联盟不断改变，才使一场被宣称结束的战事成为下一次动员的背景。对迁徙家庭而言，边地战争不是抽象的军费，而是去向、编户、土地和保护者都可能改变的过程。

\subsection{州郡为什么有了兵}

一八四年黄巾起事与凉州军民、羌胡部众的动乱并行。广宗、下曲阳等地的主力受挫不等于乡里、佃作和逃户立即恢复。为应对多线战争，朝廷在一八八年扩大州牧、刺史的军政用途；制度调整可靠，各州实际兵额、财力和权力却很不相同。不能把后来的割据当成当年唯一的设计，但可以看到中央已不得不借州郡将领、豪强和边军解决本应由共同军队承担的压力。

幽州的张纯、张举之乱、乌桓关系的变化，以及凉州军人网络，都让地方将领拥有自筹、招募与谈判的空间。董卓进入雒阳不是一个将领凭野心从边地闯入都城，而是长期征发所形成的军政资源在宫廷危机中取得入口。洛阳试图以西园八校尉重建直接武装，也未能把彼此竞争的兵源重新置于同一指挥。

\eventanchor{EH-0177-XIANBEI-RAIDS}{熹平六年对鲜卑的出击受挫}，进一步暴露了固定防线与远程追击之间的供给难题。边患及出师可以确定，檀石槐的政治统一程度、军数和损失却不能被后来的叙事磨平。对北边居民、州郡和将领来说，防务不是朝廷在地图边缘画的一条线，而是牧地、道路、马匹和能否取得粮饷的连续协商；正是在这些协商中，地方筹措武装的理由愈来愈强。

\subsection{许县的粮与皇帝}

曹操在一九六年迎献帝到许县，把皇帝诏令、屯田、军队与行政人员重新编入北方的军政框架。屯田的起始、规模和地区差异仍有争议，不能写成无阻执行的万能制度；它至少说明，仅有皇帝名义不足以维持竞争，任何集团都必须持续解决人和粮。二〇七年北征乌桓后，河北、幽州道路更深地进入曹氏控制，汉廷剩余的名义与北方资源的实际占有也更清楚地分离。

\begin{turningpoint}[地方化不是机构突然消失]
州郡、边将与豪强取得军权，并不表示中央机构在某日消散。征发、任命、道路、粮仓和皇帝名义被不同集团分别接管，原来连接帝国的接口开始服务于彼此竞争的中心。
\end{turningpoint}

\begin{sourcenote}[边郡成本的证据]
《后汉书·西羌传》、相关帝纪、段颎传及《后汉书·董卓传》保存战事和朝廷处置；它们不能替代连续军粮账、流民名册或各部内部意见。黄巾、凉州与州牧扩权的年序可相互校对，实际覆盖范围和社会后果必须保留限度。
\end{sourcenote}
\subsection{从凉州到许县：军费如何留下政治后果}

凉州的迁徙、招募和自筹粮饷，与许县把屯田、皇帝和北方军队重新接合，并不是同一条由边地直达中枢的因果线；它们却共同显示，持续供给一旦不能在共同规则内完成，能够筹兵筹粮的人便会获得新的议价位置。《后汉书·灵帝纪》《后汉书·董卓传》可追踪危机中的朝廷处置，《三国志·武帝纪》《后汉书·献帝纪》与《资治通鉴》可核对许县后的年序。它们记录显著行动和官署语言，不能据此量化基层负担、兵数或私下动机。
\subsection{账本增补的编年桥}
这些首次出现的锚点以账本的纪年、行动者与史料限度串联；它们只承载可核的桥接，不把简略记载延伸为未经证实的场景。

\noindent 永平十六年（73年），东汉与北匈奴面对北匈奴仍威胁河西交通，东汉试图以伊吾卢和绿洲盟友恢复通往西域的信息与补给路线；\eventanchor{EH-0073-DOU-GU-YIWULU}{窦固等出击北匈奴，取伊吾卢；班超奉使鄯善，东汉重启西域经营}，其后汉廷重获进入西域的军事据点，但绿洲政权的服属需要持续外交和驻军，非一次远征所能固定。记事依据《后汉书·明帝纪》；《后汉书·窦固传》；《后汉书·班超传》；出师年份可靠；各路战果与班超使节行动的先后有争议。
\par

\noindent 永平十八年（75年）的东汉、西域诸国与北匈奴发生\eventanchor{EH-0075-CHEN-MU-CRISIS}{北匈奴攻车师，西域都护陈睦遇害，东汉一度撤出西域据点}。东汉前推据点依赖河西补给和绿洲合作，北匈奴可利用车师地缘切断驻军联系构成其条件，东汉西域经营暂时收缩，显示都护名义与可持续的交通、粮秣和盟友并不等同随之进入后续年序。《后汉书·明帝纪》；《后汉书·耿恭传》；《后汉书·西域传》提供这一线索，但核心危机可靠；撤军范围与北匈奴、车师各部责任有争议。
\par

\noindent 在建初元年（76年），\eventanchor{EH-0076-QIANG-REBELLION}{西羌在陇西、金城等地反叛，东汉长期羌乱由此加剧}把东汉与西羌置于边郡官吏处置、徭役与牧地冲突叠加，羌部内部联盟又随战争重组与陇右驻军、军费和迁徙压力持续增加；史书族名不应被视为固定而统一的社会实体之间。可据《后汉书·西羌传》；《后汉书·章帝纪》；《后汉纪》卷三十七追踪其大体经过；起事年份大体可靠；部落称谓、参与范围和初期战事有争议。
\par

\noindent 东汉与北匈奴在永元元年（89年）的\eventanchor{EH-0089-DOU-XIAN-JILU}{窦宪等北出稽落山，大破北匈奴，立燕然山铭纪功}，并非孤立一事：窦氏希望以大规模北伐巩固辅政威望，南匈奴与边郡骑兵为行动提供协力，且北匈奴受到重创，窦宪军功显著抬升；铭文是胜利宣示，不能直接作为战果总数的独立证明。材料见《后汉书·和帝纪》；《后汉书·窦宪传》；《后汉书·南匈奴传》；燕然山铭；战役年月可靠；兵数、斩获与铭文功绩措辞有夸饰可能。
\par

\noindent \eventanchor{EH-0091-BAN-CHAO-PROTECTOR}{班超受命为西域都护，东汉在多处绿洲重新建立使节与军事网络}发生于永元三年（91年），行动者为东汉与西域诸国。它回应北匈奴势力受挫后，东汉试图将临时使团关系转化为较稳定的都护协调，并把汉廷恢复对部分绿洲的影响，但该体系依赖地方王权、商路和补给，不能表述为完整的直接统治留给随后的人群和官署。《后汉书·班超传》；《后汉书·西域传》；《后汉纪》卷五十二可作核验，惟任命年份可靠；各国服属程度和驻军规模有争议。
\par

\noindent 从东汉西域网络暂获恢复，班超希望获得更远西方政权和海陆路线的情报出发，东汉、安息与西域诸国于永元九年（97年）出现\eventanchor{EH-0097-GAN-YING-ANXI}{班超遣甘英西使，至安息属地而未能到达大秦}；汉文史料保存了经安息获得的西方地理知识，但旅行未达终点，不能据此推定汉罗直接外交关系。这里采用《后汉书·西域传》；《后汉书·班超传》；《后汉纪》卷五十八的记载范围，不越过遣使事实可靠；到达地点、海路信息和“大秦”识别有争议。
\par

\noindent 永元十四年（102年），东汉与西域诸国的\eventanchor{EH-0102-BAN-CHAO-DEATH}{班超返洛阳后去世，东汉西域经营失去长期协调者}使班超长期依靠地方盟友、军队和个人威望维系都护关系，年老请求返京转为可见的行动，西域体系的脆弱性更加显现，后任难以完全接续其个人网络，东汉投入与收益持续受质疑。其依据为《后汉书·班超传》；《后汉书·和帝纪》；《后汉纪》卷六十三，归还与卒年可靠；其离任对各绿洲政治的直接影响有争议。
\par

\noindent \eventanchor{EH-0167-DUAN-JIONG-QIANG}{段颎在陇右连续进军，东羌大部被击散或降附，战事伴随严酷杀戮}这一永康元年（167年）至建宁二年（169年）的节点，关联东汉与东羌、长期羌乱消耗财政并威胁陇右交通，朝廷授权段颎实施高强度围剿和短期内恢复部分通道与郡县控制，却加重地方仇怨和军政依赖，后续凉州危机并未根除。《后汉书·段颎传》；《后汉书·西羌传》；《后汉纪》卷九十三至九十五记录了它的基本轮廓；战役序列大体可靠；斩获、迁徙和“平定”范围有争议。
\par
\noindent 建武二十四年（48年），东汉与南匈奴面临的不是抽象的变局：匈奴单于继承与资源分配冲突使比部寻求外援，东汉希望以附属部众缓冲北边。这构成南匈奴日逐王比归附汉廷，南北匈奴分裂，汉在边郡安置其部众的条件，却不预定行动将如何发生。 \eventanchor{EH-LEDGER-012-1}{南匈奴日逐王比归附汉廷，南北匈奴分裂，汉在边郡安置其部众：压力与起点}。
\par
\noindent 就这次行动而言，《后汉书·南匈奴传》；《后汉纪》卷二十一；《资治通鉴》卷四十四提供了最小的年序骨架：南匈奴日逐王比归附汉廷，南北匈奴分裂，汉在边郡安置其部众。它们可支撑事件存在与大体先后，不能补齐未被记录的环节。 \eventanchor{EH-LEDGER-012-2}{南匈奴日逐王比归附汉廷，南北匈奴分裂，汉在边郡安置其部众：进入年序的行动}。
\par
\noindent 这一节点带来的直接变化是：南匈奴成为汉廷边防伙伴也保有自身首领网络，汉匈关系转为分裂政权间的长期协商。它改变的是后来行动可利用的资源和位置，未必立刻改造每一处地方社会。 \eventanchor{EH-LEDGER-012-3}{南匈奴日逐王比归附汉廷，南北匈奴分裂，汉在边郡安置其部众：可见的转折}。
\par
\noindent 关于南匈奴日逐王比归附汉廷，南北匈奴分裂，汉在边郡安置其部众，《后汉书·南匈奴传》；《后汉纪》卷二十一；《资治通鉴》卷四十四提供可检的后出史家叙事依据。由此可以约束时间与主体，却不能把文本的沉默填成具体场景。 \eventanchor{EH-LEDGER-012-4}{南匈奴日逐王比归附汉廷，南北匈奴分裂，汉在边郡安置其部众：空间与执行}。
\par
\noindent 现代研究可沿“南匈奴内附、边郡安置与草原政治研究”继续追问南匈奴日逐王比归附汉廷，南北匈奴分裂，汉在边郡安置其部众，以免把传世叙事中的概括误作已经完成的解释。 \eventanchor{EH-LEDGER-012-5}{南匈奴日逐王比归附汉廷，南北匈奴分裂，汉在边郡安置其部众：后续压力}。
\par
\noindent 本节点的可说范围止于“归附年份可靠；部众数量、驻牧边界和自主程度有争议”。在这个界限内，南匈奴日逐王比归附汉廷，南北匈奴分裂，汉在边郡安置其部众连接前后年序；越过它便是无据补写。 \eventanchor{EH-LEDGER-012-6}{南匈奴日逐王比归附汉廷，南北匈奴分裂，汉在边郡安置其部众：材料限度}。
\par
\noindent 东汉与北匈奴在永平十六年（73年）所要处理的是北匈奴仍威胁河西交通，东汉试图以伊吾卢和绿洲盟友恢复通往西域的信息与补给路线；窦固等出击北匈奴，取伊吾卢；班超奉使鄯善，东汉重启西域经营由此成为可见的节点，而非一条自动展开的结果。 \eventanchor{EH-LEDGER-015-1}{窦固等出击北匈奴，取伊吾卢；班超奉使鄯善，东汉重启西域经营：压力与起点}。
\par
\noindent 年序所能抓住的动作是窦固等出击北匈奴，取伊吾卢；班超奉使鄯善，东汉重启西域经营。《后汉书·明帝纪》；《后汉书·窦固传》；《后汉书·班超传》把它置入前后事件之间；叙事足以确认这一节点，不足以替每一位参与者写出完整经历。 \eventanchor{EH-LEDGER-015-2}{窦固等出击北匈奴，取伊吾卢；班超奉使鄯善，东汉重启西域经营：进入年序的行动}。
\par
\noindent 因此，窦固等出击北匈奴，取伊吾卢；班超奉使鄯善，东汉重启西域经营之后出现了汉廷重获进入西域的军事据点，但绿洲政权的服属需要持续外交和驻军，非一次远征所能固定。这一结果说明条件被重新排列，不表示所有行动者以同一方式接受它。 \eventanchor{EH-LEDGER-015-3}{窦固等出击北匈奴，取伊吾卢；班超奉使鄯善，东汉重启西域经营：可见的转折}。
\par
\noindent 《后汉书·明帝纪》；《后汉书·窦固传》；《后汉书·班超传》保存的是窦固等出击北匈奴，取伊吾卢；班超奉使鄯善，东汉重启西域经营的后出史家叙事材料线索。它使日期、主体或行动得以定位；空间尺度与实际执行须另作核验。 \eventanchor{EH-LEDGER-015-4}{窦固等出击北匈奴，取伊吾卢；班超奉使鄯善，东汉重启西域经营：空间与执行}。
\par
\noindent 关于窦固等出击北匈奴，取伊吾卢；班超奉使鄯善，东汉重启西域经营，可进一步参照永平北伐、伊吾卢与西域使团研究。研究问题能揭示盲点，却不应将推测写回为确定事实。 \eventanchor{EH-LEDGER-015-5}{窦固等出击北匈奴，取伊吾卢；班超奉使鄯善，东汉重启西域经营：后续压力}。
\par
\noindent 证据边界在于：出师年份可靠；各路战果与班超使节行动的先后有争议。因此本段只把窦固等出击北匈奴，取伊吾卢；班超奉使鄯善，东汉重启西域经营安放在可核的年序、行动者和后果之间。 \eventanchor{EH-LEDGER-015-6}{窦固等出击北匈奴，取伊吾卢；班超奉使鄯善，东汉重启西域经营：材料限度}。
\par
\noindent 在永平十八年（75年），北匈奴攻车师，西域都护陈睦遇害，东汉一度撤出西域据点之前已有一层具体压力：东汉前推据点依赖河西补给和绿洲合作，北匈奴可利用车师地缘切断驻军联系。它让东汉、西域诸国与北匈奴不得不在既有资源与信息中作出选择。 \eventanchor{EH-LEDGER-016-1}{北匈奴攻车师，西域都护陈睦遇害，东汉一度撤出西域据点：压力与起点}。
\par
\noindent 《后汉书·明帝纪》；《后汉书·耿恭传》；《后汉书·西域传》记录了北匈奴攻车师，西域都护陈睦遇害，东汉一度撤出西域据点这一转折。它们让永平十八年（75年）的基本顺序可以核对，人物言辞、具体人数或内部商议仍须另证。 \eventanchor{EH-LEDGER-016-2}{北匈奴攻车师，西域都护陈睦遇害，东汉一度撤出西域据点：进入年序的行动}。
\par
\noindent 记录所允许追出的下一步是：东汉西域经营暂时收缩，显示都护名义与可持续的交通、粮秣和盟友并不等同。后续变化需在不同地区分别检验，不能由这一个节点概括全国。 \eventanchor{EH-LEDGER-016-3}{北匈奴攻车师，西域都护陈睦遇害，东汉一度撤出西域据点：可见的转折}。
\par
\noindent 要把北匈奴攻车师，西域都护陈睦遇害，东汉一度撤出西域据点落在可核材料上，应回到《后汉书·明帝纪》；《后汉书·耿恭传》；《后汉书·西域传》。这些后出史家叙事材料能限定叙述的骨架，不能凭自身复原所有地点与人员。 \eventanchor{EH-LEDGER-016-4}{北匈奴攻车师，西域都护陈睦遇害，东汉一度撤出西域据点：空间与执行}。
\par
\noindent 本卷把北匈奴攻车师，西域都护陈睦遇害，东汉一度撤出西域据点留给“车师危机、都护制度与西域驻军补给研究”这一路问题意识继续检验；解释的分歧不被伪装为材料的一致。 \eventanchor{EH-LEDGER-016-5}{北匈奴攻车师，西域都护陈睦遇害，东汉一度撤出西域据点：后续压力}。
\par
\noindent 核心危机可靠；撤军范围与北匈奴、车师各部责任有争议。这也是为何北匈奴攻车师，西域都护陈睦遇害，东汉一度撤出西域据点在此只作来源感知的首次描写，不把不稳的细节写成结论。 \eventanchor{EH-LEDGER-016-6}{北匈奴攻车师，西域都护陈睦遇害，东汉一度撤出西域据点：材料限度}。
\par
\noindent 建初元年（76年），东汉与西羌面临的不是抽象的变局：边郡官吏处置、徭役与牧地冲突叠加，羌部内部联盟又随战争重组。这构成西羌在陇西、金城等地反叛，东汉长期羌乱由此加剧的条件，却不预定行动将如何发生。 \eventanchor{EH-LEDGER-018-1}{西羌在陇西、金城等地反叛，东汉长期羌乱由此加剧：压力与起点}。
\par
\noindent 就这次行动而言，《后汉书·西羌传》；《后汉书·章帝纪》；《后汉纪》卷三十七提供了最小的年序骨架：西羌在陇西、金城等地反叛，东汉长期羌乱由此加剧。它们可支撑事件存在与大体先后，不能补齐未被记录的环节。 \eventanchor{EH-LEDGER-018-2}{西羌在陇西、金城等地反叛，东汉长期羌乱由此加剧：进入年序的行动}。
\par
\noindent 这一节点带来的直接变化是：陇右驻军、军费和迁徙压力持续增加；史书族名不应被视为固定而统一的社会实体。它改变的是后来行动可利用的资源和位置，未必立刻改造每一处地方社会。 \eventanchor{EH-LEDGER-018-3}{西羌在陇西、金城等地反叛，东汉长期羌乱由此加剧：可见的转折}。
\par
\noindent 关于西羌在陇西、金城等地反叛，东汉长期羌乱由此加剧，《后汉书·西羌传》；《后汉书·章帝纪》；《后汉纪》卷三十七提供可检的后出史家叙事依据。由此可以约束时间与主体，却不能把文本的沉默填成具体场景。 \eventanchor{EH-LEDGER-018-4}{西羌在陇西、金城等地反叛，东汉长期羌乱由此加剧：空间与执行}。
\par
\noindent 现代研究可沿“东汉羌乱、边郡征发与族群政治研究”继续追问西羌在陇西、金城等地反叛，东汉长期羌乱由此加剧，以免把传世叙事中的概括误作已经完成的解释。 \eventanchor{EH-LEDGER-018-5}{西羌在陇西、金城等地反叛，东汉长期羌乱由此加剧：后续压力}。
\par
\noindent 本节点的可说范围止于“起事年份大体可靠；部落称谓、参与范围和初期战事有争议”。在这个界限内，西羌在陇西、金城等地反叛，东汉长期羌乱由此加剧连接前后年序；越过它便是无据补写。 \eventanchor{EH-LEDGER-018-6}{西羌在陇西、金城等地反叛，东汉长期羌乱由此加剧：材料限度}。
\par
\noindent 东汉与北匈奴在永元元年（89年）所要处理的是窦氏希望以大规模北伐巩固辅政威望，南匈奴与边郡骑兵为行动提供协力；窦宪等北出稽落山，大破北匈奴，立燕然山铭纪功由此成为可见的节点，而非一条自动展开的结果。 \eventanchor{EH-LEDGER-021-1}{窦宪等北出稽落山，大破北匈奴，立燕然山铭纪功：压力与起点}。
\par
\noindent 年序所能抓住的动作是窦宪等北出稽落山，大破北匈奴，立燕然山铭纪功。《后汉书·和帝纪》；《后汉书·窦宪传》；《后汉书·南匈奴传》；燕然山铭把它置入前后事件之间；叙事足以确认这一节点，不足以替每一位参与者写出完整经历。 \eventanchor{EH-LEDGER-021-2}{窦宪等北出稽落山，大破北匈奴，立燕然山铭纪功：进入年序的行动}。
\par
\noindent 因此，窦宪等北出稽落山，大破北匈奴，立燕然山铭纪功之后出现了北匈奴受到重创，窦宪军功显著抬升；铭文是胜利宣示，不能直接作为战果总数的独立证明。这一结果说明条件被重新排列，不表示所有行动者以同一方式接受它。 \eventanchor{EH-LEDGER-021-3}{窦宪等北出稽落山，大破北匈奴，立燕然山铭纪功：可见的转折}。
\par
\noindent 《后汉书·和帝纪》；《后汉书·窦宪传》；《后汉书·南匈奴传》；燕然山铭保存的是窦宪等北出稽落山，大破北匈奴，立燕然山铭纪功的出土文字或铭刻材料线索。它使日期、主体或行动得以定位；空间尺度与实际执行须另作核验。 \eventanchor{EH-LEDGER-021-4}{窦宪等北出稽落山，大破北匈奴，立燕然山铭纪功：空间与执行}。
\par
\noindent 关于窦宪等北出稽落山，大破北匈奴，立燕然山铭纪功，可进一步参照窦宪北伐、燕然山铭与东汉边疆纪功研究。研究问题能揭示盲点，却不应将推测写回为确定事实。 \eventanchor{EH-LEDGER-021-5}{窦宪等北出稽落山，大破北匈奴，立燕然山铭纪功：后续压力}。
\par
\noindent 证据边界在于：战役年月可靠；兵数、斩获与铭文功绩措辞有夸饰可能。因此本段只把窦宪等北出稽落山，大破北匈奴，立燕然山铭纪功安放在可核的年序、行动者和后果之间。 \eventanchor{EH-LEDGER-021-6}{窦宪等北出稽落山，大破北匈奴，立燕然山铭纪功：材料限度}。
\par
\noindent 在永元三年（91年），班超受命为西域都护，东汉在多处绿洲重新建立使节与军事网络之前已有一层具体压力：北匈奴势力受挫后，东汉试图将临时使团关系转化为较稳定的都护协调。它让东汉与西域诸国不得不在既有资源与信息中作出选择。 \eventanchor{EH-LEDGER-022-1}{班超受命为西域都护，东汉在多处绿洲重新建立使节与军事网络：压力与起点}。
\par
\noindent 《后汉书·班超传》；《后汉书·西域传》；《后汉纪》卷五十二记录了班超受命为西域都护，东汉在多处绿洲重新建立使节与军事网络这一转折。它们让永元三年（91年）的基本顺序可以核对，人物言辞、具体人数或内部商议仍须另证。 \eventanchor{EH-LEDGER-022-2}{班超受命为西域都护，东汉在多处绿洲重新建立使节与军事网络：进入年序的行动}。
\par
\noindent 记录所允许追出的下一步是：汉廷恢复对部分绿洲的影响，但该体系依赖地方王权、商路和补给，不能表述为完整的直接统治。后续变化需在不同地区分别检验，不能由这一个节点概括全国。 \eventanchor{EH-LEDGER-022-3}{班超受命为西域都护，东汉在多处绿洲重新建立使节与军事网络：可见的转折}。
\par
\noindent 要把班超受命为西域都护，东汉在多处绿洲重新建立使节与军事网络落在可核材料上，应回到《后汉书·班超传》；《后汉书·西域传》；《后汉纪》卷五十二。这些后出史家叙事材料能限定叙述的骨架，不能凭自身复原所有地点与人员。 \eventanchor{EH-LEDGER-022-4}{班超受命为西域都护，东汉在多处绿洲重新建立使节与军事网络：空间与执行}。
\par
\noindent 本卷把班超受命为西域都护，东汉在多处绿洲重新建立使节与军事网络留给“班超都护、西域绿洲政治与汉军驻防研究”这一路问题意识继续检验；解释的分歧不被伪装为材料的一致。 \eventanchor{EH-LEDGER-022-5}{班超受命为西域都护，东汉在多处绿洲重新建立使节与军事网络：后续压力}。
\par
\noindent 任命年份可靠；各国服属程度和驻军规模有争议。这也是为何班超受命为西域都护，东汉在多处绿洲重新建立使节与军事网络在此只作来源感知的首次描写，不把不稳的细节写成结论。 \eventanchor{EH-LEDGER-022-6}{班超受命为西域都护，东汉在多处绿洲重新建立使节与军事网络：材料限度}。
\par
\noindent 永元九年（97年），东汉、安息与西域诸国面临的不是抽象的变局：东汉西域网络暂获恢复，班超希望获得更远西方政权和海陆路线的情报。这构成班超遣甘英西使，至安息属地而未能到达大秦的条件，却不预定行动将如何发生。 \eventanchor{EH-LEDGER-024-1}{班超遣甘英西使，至安息属地而未能到达大秦：压力与起点}。
\par
\noindent 就这次行动而言，《后汉书·西域传》；《后汉书·班超传》；《后汉纪》卷五十八提供了最小的年序骨架：班超遣甘英西使，至安息属地而未能到达大秦。它们可支撑事件存在与大体先后，不能补齐未被记录的环节。 \eventanchor{EH-LEDGER-024-2}{班超遣甘英西使，至安息属地而未能到达大秦：进入年序的行动}。
\par
\noindent 这一节点带来的直接变化是：汉文史料保存了经安息获得的西方地理知识，但旅行未达终点，不能据此推定汉罗直接外交关系。它改变的是后来行动可利用的资源和位置，未必立刻改造每一处地方社会。 \eventanchor{EH-LEDGER-024-3}{班超遣甘英西使，至安息属地而未能到达大秦：可见的转折}。
\par
\noindent 关于班超遣甘英西使，至安息属地而未能到达大秦，《后汉书·西域传》；《后汉书·班超传》；《后汉纪》卷五十八提供可检的后出史家叙事依据。由此可以约束时间与主体，却不能把文本的沉默填成具体场景。 \eventanchor{EH-LEDGER-024-4}{班超遣甘英西使，至安息属地而未能到达大秦：空间与执行}。
\par
\noindent 现代研究可沿“甘英西使、安息中介与欧亚交通信息研究”继续追问班超遣甘英西使，至安息属地而未能到达大秦，以免把传世叙事中的概括误作已经完成的解释。 \eventanchor{EH-LEDGER-024-5}{班超遣甘英西使，至安息属地而未能到达大秦：后续压力}。
\par
\noindent 本节点的可说范围止于“遣使事实可靠；到达地点、海路信息和“大秦”识别有争议”。在这个界限内，班超遣甘英西使，至安息属地而未能到达大秦连接前后年序；越过它便是无据补写。 \eventanchor{EH-LEDGER-024-6}{班超遣甘英西使，至安息属地而未能到达大秦：材料限度}。
\par
\noindent 把班超返洛阳后去世，东汉西域经营失去长期协调者放回永元十四年（102年），先要看东汉与西域诸国所处的约束：班超长期依靠地方盟友、军队和个人威望维系都护关系，年老请求返京。史料可追到这种压力，不能把它简化成唯一动机。 \eventanchor{EH-LEDGER-025-1}{班超返洛阳后去世，东汉西域经营失去长期协调者：压力与起点}。
\par
\noindent 永元十四年（102年）的叙事在班超返洛阳后去世，东汉西域经营失去长期协调者处汇合。依据是《后汉书·班超传》；《后汉书·和帝纪》；《后汉纪》卷六十三；可核的是该节点及其时序，细部仍要服从各书的立场与成书距离。 \eventanchor{EH-LEDGER-025-2}{班超返洛阳后去世，东汉西域经营失去长期协调者：进入年序的行动}。
\par
\noindent 其后果首先落在可接续的关系上：西域体系的脆弱性更加显现，后任难以完全接续其个人网络，东汉投入与收益持续受质疑。从这一变化到更远的制度或社会影响，仍隔着地方执行与时间。 \eventanchor{EH-LEDGER-025-3}{班超返洛阳后去世，东汉西域经营失去长期协调者：可见的转折}。
\par
\noindent 本段把班超返洛阳后去世，东汉西域经营失去长期协调者系在《后汉书·班超传》；《后汉书·和帝纪》；《后汉纪》卷六十三这一材料路径上。作为后出史家叙事，它的证明力随所记事项而变，不能外推为全面记录。 \eventanchor{EH-LEDGER-025-4}{班超返洛阳后去世，东汉西域经营失去长期协调者：空间与执行}。
\par
\noindent “班超身后西域、个人网络与都护制度研究”是理解班超返洛阳后去世，东汉西域经营失去长期协调者的一条研究路径；它帮助追问机制，不替代本段所依据的史料。 \eventanchor{EH-LEDGER-025-5}{班超返洛阳后去世，东汉西域经营失去长期协调者：后续压力}。
\par
\noindent 归还与卒年可靠；其离任对各绿洲政治的直接影响有争议；这限制了班超返洛阳后去世，东汉西域经营失去长期协调者可被叙述的密度，也保留了史料没有回答的问题。 \eventanchor{EH-LEDGER-025-6}{班超返洛阳后去世，东汉西域经营失去长期协调者：材料限度}。
\par
\noindent 在永初元年（107年），西羌在凉州、并州一带再起，东汉长期投入边郡军费与移民安置之前已有一层具体压力：官府征发、边郡冲突和羌部联盟变化使此前平叛未能消除结构性矛盾。它让东汉与西羌不得不在既有资源与信息中作出选择。 \eventanchor{EH-LEDGER-028-1}{西羌在凉州、并州一带再起，东汉长期投入边郡军费与移民安置：压力与起点}。
\par
\noindent 《后汉书·安帝纪》；《后汉书·西羌传》；《后汉纪》卷六十八记录了西羌在凉州、并州一带再起，东汉长期投入边郡军费与移民安置这一转折。它们让永初元年（107年）的基本顺序可以核对，人物言辞、具体人数或内部商议仍须另证。 \eventanchor{EH-LEDGER-028-2}{西羌在凉州、并州一带再起，东汉长期投入边郡军费与移民安置：进入年序的行动}。
\par
\noindent 记录所允许追出的下一步是：凉州成为持续军事财政负担，流民、降附与迁徙改变地方社会，正史军功数字须谨慎使用。后续变化需在不同地区分别检验，不能由这一个节点概括全国。 \eventanchor{EH-LEDGER-028-3}{西羌在凉州、并州一带再起，东汉长期投入边郡军费与移民安置：可见的转折}。
\par
\noindent 要把西羌在凉州、并州一带再起，东汉长期投入边郡军费与移民安置落在可核材料上，应回到《后汉书·安帝纪》；《后汉书·西羌传》；《后汉纪》卷六十八。这些后出史家叙事材料能限定叙述的骨架，不能凭自身复原所有地点与人员。 \eventanchor{EH-LEDGER-028-4}{西羌在凉州、并州一带再起，东汉长期投入边郡军费与移民安置：空间与执行}。
\par
\noindent 本卷把西羌在凉州、并州一带再起，东汉长期投入边郡军费与移民安置留给“永初羌乱、凉州财政与边民迁徙研究”这一路问题意识继续检验；解释的分歧不被伪装为材料的一致。 \eventanchor{EH-LEDGER-028-5}{西羌在凉州、并州一带再起，东汉长期投入边郡军费与移民安置：后续压力}。
\par
\noindent 战乱开始年份可靠；参战部落和各郡受害程度有争议。这也是为何西羌在凉州、并州一带再起，东汉长期投入边郡军费与移民安置在此只作来源感知的首次描写，不把不稳的细节写成结论。 \eventanchor{EH-LEDGER-028-6}{西羌在凉州、并州一带再起，东汉长期投入边郡军费与移民安置：材料限度}。
\par

\noindent 在党锢列传：熹平五年，\eventanchor{EH-DENSE-0221}{熹平五年，永昌太守曹鸾上书大讼党人，言甚方切}把东汉与边地、地方社会置于同卷前后条目所示的制度、地方或边疆条件与作为跨材料核验的编年节点保留之间。可据《后汉书》·党锢列传追踪其大体经过；《后汉书》传志兼具官修分类与后出叙事；地方实践、数字和群体名称须与简牍、碑刻及考古交叉核验。
\par

\noindent 东汉与边地、地方社会在循吏列传：建武二年的\eventanchor{EH-DENSE-0222}{飒下车，修庠序之教，设婚姻之礼}，并非孤立一事：同卷前后条目所示的制度、地方或边疆条件，且作为跨材料核验的编年节点保留。材料见《后汉书》·循吏列传；《后汉书》传志兼具官修分类与后出叙事；地方实践、数字和群体名称须与简牍、碑刻及考古交叉核验。
\par

\noindent \eventanchor{EH-DENSE-0223}{帝念觽功美，封为鄛乡侯，食邑千五百户}发生于宦者列传：十四年，行动者为东汉与边地、地方社会。它回应同卷前后条目所示的制度、地方或边疆条件，并把作为跨材料核验的编年节点保留留给随后的人群和官署。《后汉书》·宦者列传可作核验，惟《后汉书》传志兼具官修分类与后出叙事；地方实践、数字和群体名称须与简牍、碑刻及考古交叉核验。
\par

\noindent 从同卷前后条目所示的制度、地方或边疆条件出发，东汉与边地、地方社会于儒林列传：建武五年出现\eventanchor{EH-DENSE-0224}{建武五年，乃修起太学，稽式古典，笾豆干戚之容，备之于列，服方领习矩步者，委它乎其中}；作为跨材料核验的编年节点保留。这里采用《后汉书》·儒林列传的记载范围，不越过《后汉书》传志兼具官修分类与后出叙事；地方实践、数字和群体名称须与简牍、碑刻及考古交叉核验。
\par

\noindent 儒林列传：十二年，东汉与边地、地方社会的\eventanchor{EH-DENSE-0225}{泽大怒，以妻干犯斋禁，遂收送诏狱谢罪}使同卷前后条目所示的制度、地方或边疆条件转为可见的行动，作为跨材料核验的编年节点保留。其依据为《后汉书》·儒林列传，《后汉书》传志兼具官修分类与后出叙事；地方实践、数字和群体名称须与简牍、碑刻及考古交叉核验。
\par

\noindent \eventanchor{EH-DENSE-0226}{王莽辅政，召为郎，稍迁酒泉都尉}这一方术列传：八年的节点，关联东汉与边地、地方社会、同卷前后条目所示的制度、地方或边疆条件和作为跨材料核验的编年节点保留。《后汉书》·方术列传记录了它的基本轮廓；《后汉书》传志兼具官修分类与后出叙事；地方实践、数字和群体名称须与简牍、碑刻及考古交叉核验。
\par

\noindent 东汉与边地、地方社会在方术列传：元初七年处理同卷前后条目所示的制度、地方或边疆条件时，出现\eventanchor{EH-DENSE-0227}{元初七年，郡界有芝草生，太守刘祗欲上言之，以问□}；作为跨材料核验的编年节点保留。可核材料为《后汉书》·方术列传，且《后汉书》传志兼具官修分类与后出叙事；地方实践、数字和群体名称须与简牍、碑刻及考古交叉核验。
\par

\noindent 列女传：永平三年的\eventanchor{EH-DENSE-0228}{我言而不用，君必谓我不奉教令，则罪在我矣}，让东汉与边地、地方社会的同卷前后条目所示的制度、地方或边疆条件在具体事务中显形；作为跨材料核验的编年节点保留。《后汉书》·列女传所能支持的范围仍受《后汉书》传志兼具官修分类与后出叙事；地方实践、数字和群体名称须与简牍、碑刻及考古交叉核验限制。
\par

\noindent 对东汉与边地、地方社会而言，\eventanchor{EH-DENSE-0229}{建武中元二年，倭奴国奉贡朝贺，使人自称大夫，倭国之极南界也}发生在东夷列传：建武中元二年：同卷前后条目所示的制度、地方或边疆条件，随后作为跨材料核验的编年节点保留。参看《后汉书》·东夷列传；《后汉书》传志兼具官修分类与后出叙事；地方实践、数字和群体名称须与简牍、碑刻及考古交叉核验。
\par

\noindent 同卷前后条目所示的制度、地方或边疆条件之下，东汉与边地、地方社会于南蛮西南夷列传：和帝永元四年冬经历\eventanchor{EH-DENSE-0230}{和帝永元四年冬，溇中、澧中蛮潭戎等反，燔烧邮亭，杀略吏民，郡兵击破降之}，并面对作为跨材料核验的编年节点保留。此处依《后汉书》·南蛮西南夷列传，《后汉书》传志兼具官修分类与后出叙事；地方实践、数字和群体名称须与简牍、碑刻及考古交叉核验。
\par

\noindent 西羌传：永平元年，东汉与边地、地方社会面对同卷前后条目所示的制度、地方或边疆条件；\eventanchor{EH-DENSE-0231}{永平元年，复遣中郎将窦固、捕虏将军马武等击滇吾于西邯，大破之}，其后作为跨材料核验的编年节点保留。记事依据《后汉书》·西羌传；《后汉书》传志兼具官修分类与后出叙事；地方实践、数字和群体名称须与简牍、碑刻及考古交叉核验。
\par

\noindent 西域传：九年的东汉与边地、地方社会发生\eventanchor{EH-DENSE-0232}{于是远国蒙奇、兜勒皆来归服，遣使贡献}。同卷前后条目所示的制度、地方或边疆条件构成其条件，作为跨材料核验的编年节点保留随之进入后续年序。《后汉书》·西域传提供这一线索，但《后汉书》传志兼具官修分类与后出叙事；地方实践、数字和群体名称须与简牍、碑刻及考古交叉核验。
\par

\noindent 在南匈奴列传：二十四年，\eventanchor{EH-DENSE-0233}{八部大人共议立比为呼韩邪单于，以其大父尝依汉得安，故欲袭其号}把东汉与边地、地方社会置于同卷前后条目所示的制度、地方或边疆条件与作为跨材料核验的编年节点保留之间。可据《后汉书》·南匈奴列传追踪其大体经过；《后汉书》传志兼具官修分类与后出叙事；地方实践、数字和群体名称须与简牍、碑刻及考古交叉核验。
\par

\noindent 东汉与边地、地方社会在乌桓鲜卑列传：建武二十一年的\eventanchor{EH-DENSE-0234}{建武二十一年，遣伏波将军马援将三千骑出五阮关掩击之}，并非孤立一事：同卷前后条目所示的制度、地方或边疆条件，且作为跨材料核验的编年节点保留。材料见《后汉书》·乌桓鲜卑列传；《后汉书》传志兼具官修分类与后出叙事；地方实践、数字和群体名称须与简牍、碑刻及考古交叉核验。
\par

\noindent \eventanchor{EH-DENSE-0235}{帝省奏大怒，即诏司隶、益州槛车收鸾，送槐里狱掠杀之}发生于党锢列传：熹平五年，行动者为东汉与边地、地方社会。它回应同卷前后条目所示的制度、地方或边疆条件，并把作为跨材料核验的编年节点保留留给随后的人群和官署。《后汉书》·党锢列传可作核验，惟《后汉书》传志兼具官修分类与后出叙事；地方实践、数字和群体名称须与简牍、碑刻及考古交叉核验。
\par

\noindent 从同卷前后条目所示的制度、地方或边疆条件出发，东汉与边地、地方社会于循吏列传：建武二年出现\eventanchor{EH-DENSE-0236}{民居深山，滨溪谷，习其风土，不出田租}；作为跨材料核验的编年节点保留。这里采用《后汉书》·循吏列传的记载范围，不越过《后汉书》传志兼具官修分类与后出叙事；地方实践、数字和群体名称须与简牍、碑刻及考古交叉核验。
\par

\noindent 宦者列传：元初元年，东汉与边地、地方社会的\eventanchor{EH-DENSE-0237}{桓帝延熹二年，绍封觽曾孙石雠为关内侯}使同卷前后条目所示的制度、地方或边疆条件转为可见的行动，作为跨材料核验的编年节点保留。其依据为《后汉书》·宦者列传，《后汉书》传志兼具官修分类与后出叙事；地方实践、数字和群体名称须与简牍、碑刻及考古交叉核验。
\par

\noindent \eventanchor{EH-DENSE-0238}{天子始冠通天，衣日月，备法物之驾，盛清道之仪，坐明堂而朝髃后，登灵台以望云物，袒割辟雍之上，尊养三老五更}这一儒林列传：建武五年的节点，关联东汉与边地、地方社会、同卷前后条目所示的制度、地方或边疆条件和作为跨材料核验的编年节点保留。《后汉书》·儒林列传记录了它的基本轮廓；《后汉书》传志兼具官修分类与后出叙事；地方实践、数字和群体名称须与简牍、碑刻及考古交叉核验。
\par

\noindent 东汉与边地、地方社会在儒林列传：十二年处理同卷前后条目所示的制度、地方或边疆条件时，出现\eventanchor{EH-DENSE-0239}{诏令定春秋章句，去其复重，以授皇太子}；作为跨材料核验的编年节点保留。可核材料为《后汉书》·儒林列传，且《后汉书》传志兼具官修分类与后出叙事；地方实践、数字和群体名称须与简牍、碑刻及考古交叉核验。
\par

\noindent 方术列传：八年的\eventanchor{EH-DENSE-0240}{汝南旧有鸿郄陂，成帝时，丞相翟方进奏毁败之}，让东汉与边地、地方社会的同卷前后条目所示的制度、地方或边疆条件在具体事务中显形；作为跨材料核验的编年节点保留。《后汉书》·方术列传所能支持的范围仍受《后汉书》传志兼具官修分类与后出叙事；地方实践、数字和群体名称须与简牍、碑刻及考古交叉核验限制。
\par

\noindent 对东汉与边地、地方社会而言，\eventanchor{EH-DENSE-0241}{檀以为京师当有兵气，其祸发于萧墙}发生在方术列传：元初七年：同卷前后条目所示的制度、地方或边疆条件，随后作为跨材料核验的编年节点保留。参看《后汉书》·方术列传；《后汉书》传志兼具官修分类与后出叙事；地方实践、数字和群体名称须与简牍、碑刻及考古交叉核验。
\par

\noindent 同卷前后条目所示的制度、地方或边疆条件之下，东汉与边地、地方社会于列女传：永平三年经历\eventanchor{EH-DENSE-0242}{兄固着汉书，其八表及天文志未及竟而卒，和帝诏昭就东观臧书阁踵而成之}，并面对作为跨材料核验的编年节点保留。此处依《后汉书》·列女传，《后汉书》传志兼具官修分类与后出叙事；地方实践、数字和群体名称须与简牍、碑刻及考古交叉核验。
\par

\noindent 东夷列传：建武中元二年，东汉与边地、地方社会面对同卷前后条目所示的制度、地方或边疆条件；\eventanchor{EH-DENSE-0243}{安帝永初元年，倭国王帅升等献生口百六十人，愿请见}，其后作为跨材料核验的编年节点保留。记事依据《后汉书》·东夷列传；《后汉书》传志兼具官修分类与后出叙事；地方实践、数字和群体名称须与简牍、碑刻及考古交叉核验。
\par

\noindent 南蛮西南夷列传：和帝永元四年冬的东汉与边地、地方社会发生\eventanchor{EH-DENSE-0244}{安帝元初二年，澧中蛮以郡县徭税失平，怀怨恨，遂结充中诸种二千余人，攻城杀长吏}。同卷前后条目所示的制度、地方或边疆条件构成其条件，作为跨材料核验的编年节点保留随之进入后续年序。《后汉书》·南蛮西南夷列传提供这一线索，但《后汉书》传志兼具官修分类与后出叙事；地方实践、数字和群体名称须与简牍、碑刻及考古交叉核验。
\par

\noindent 在西羌传：永平元年，\eventanchor{EH-DENSE-0245}{滇吾远引去，余悉散降，徙七千口置三辅}把东汉与边地、地方社会置于同卷前后条目所示的制度、地方或边疆条件与作为跨材料核验的编年节点保留之间。可据《后汉书》·西羌传追踪其大体经过；《后汉书》传志兼具官修分类与后出叙事；地方实践、数字和群体名称须与简牍、碑刻及考古交叉核验。
\par

\noindent 东汉与边地、地方社会在西域传：九年的\eventanchor{EH-DENSE-0246}{安帝永初元年，频攻围都护任尚、段禧等，朝廷以其险远，难相应赴，诏罢都护}，并非孤立一事：同卷前后条目所示的制度、地方或边疆条件，且作为跨材料核验的编年节点保留。材料见《后汉书》·西域传；《后汉书》传志兼具官修分类与后出叙事；地方实践、数字和群体名称须与简牍、碑刻及考古交叉核验。
\par

\noindent \eventanchor{EH-DENSE-0247}{遣弟左贤王莫将兵万余人击北单于弟薁鞬左贤王，生获之；又破北单于帐下，并得其觽合万余人，马七千匹、牛羊万头}发生于南匈奴列传：二十五年，行动者为东汉与边地、地方社会。它回应同卷前后条目所示的制度、地方或边疆条件，并把作为跨材料核验的编年节点保留留给随后的人群和官署。《后汉书》·南匈奴列传可作核验，惟《后汉书》传志兼具官修分类与后出叙事；地方实践、数字和群体名称须与简牍、碑刻及考古交叉核验。
\par

\noindent 从同卷前后条目所示的制度、地方或边疆条件出发，东汉与边地、地方社会于乌桓鲜卑列传：建武二十一年出现\eventanchor{EH-DENSE-0248}{乌桓复尾击援后，援遂晨夜奔归，比入塞，马死者千余匹}；作为跨材料核验的编年节点保留。这里采用《后汉书》·乌桓鲜卑列传的记载范围，不越过《后汉书》传志兼具官修分类与后出叙事；地方实践、数字和群体名称须与简牍、碑刻及考古交叉核验。
\par

\noindent 党锢列传：熹平五年，东汉与边地、地方社会的\eventanchor{EH-DENSE-0249}{于是又诏州郡更考党人门生故吏父子兄弟，其在位者，免官禁锢，爰及五属}使同卷前后条目所示的制度、地方或边疆条件转为可见的行动，作为跨材料核验的编年节点保留。其依据为《后汉书》·党锢列传，《后汉书》传志兼具官修分类与后出叙事；地方实践、数字和群体名称须与简牍、碑刻及考古交叉核验。
\par

\noindent \eventanchor{EH-DENSE-0250}{飒乃凿山信道五百余里，列亭传，置邮驿}这一循吏列传：建武二年的节点，关联东汉与边地、地方社会、同卷前后条目所示的制度、地方或边疆条件和作为跨材料核验的编年节点保留。《后汉书》·循吏列传记录了它的基本轮廓；《后汉书》传志兼具官修分类与后出叙事；地方实践、数字和群体名称须与简牍、碑刻及考古交叉核验。
\par

\noindent 东汉与边地、地方社会在宦者列传：元初元年处理同卷前后条目所示的制度、地方或边疆条件时，出现\eventanchor{EH-DENSE-0251}{元初元年，邓太后以伦久宿卫，封为龙亭侯，邑三百户}；作为跨材料核验的编年节点保留。可核材料为《后汉书》·宦者列传，且《后汉书》传志兼具官修分类与后出叙事；地方实践、数字和群体名称须与简牍、碑刻及考古交叉核验。
\par

\noindent 儒林列传：建武五年的\eventanchor{EH-DENSE-0252}{其后复为功臣子孙、四姓末属别立校舍，搜选高能以受其业，自期门羽林之士，悉令通孝经章句，匈奴亦遣子入学}，让东汉与边地、地方社会的同卷前后条目所示的制度、地方或边疆条件在具体事务中显形；作为跨材料核验的编年节点保留。《后汉书》·儒林列传所能支持的范围仍受《后汉书》传志兼具官修分类与后出叙事；地方实践、数字和群体名称须与简牍、碑刻及考古交叉核验限制。
\par

\noindent 对东汉与边地、地方社会而言，\eventanchor{EH-DENSE-0253}{”于是复封恭，而兴遂固辞不受爵，卒于官}发生在儒林列传：十二年：同卷前后条目所示的制度、地方或边疆条件，随后作为跨材料核验的编年节点保留。参看《后汉书》·儒林列传；《后汉书》传志兼具官修分类与后出叙事；地方实践、数字和群体名称须与简牍、碑刻及考古交叉核验。
\par

\noindent 同卷前后条目所示的制度、地方或边疆条件之下，东汉与边地、地方社会于方术列传：八年经历\eventanchor{EH-DENSE-0254}{建武中，太守邓晨欲修复其功，闻杨晓水脉，召与议之}，并面对作为跨材料核验的编年节点保留。此处依《后汉书》·方术列传，《后汉书》传志兼具官修分类与后出叙事；地方实践、数字和群体名称须与简牍、碑刻及考古交叉核验。
\par

\noindent 方术列传：元初七年，东汉与边地、地方社会面对同卷前后条目所示的制度、地方或边疆条件；\eventanchor{EH-DENSE-0255}{至延光四年，中黄门孙程扬兵殿省，诛皇后兄车骑将军阎显等，立济阴王为天子，果如所占}，其后作为跨材料核验的编年节点保留。记事依据《后汉书》·方术列传；《后汉书》传志兼具官修分类与后出叙事；地方实践、数字和群体名称须与简牍、碑刻及考古交叉核验。
\par

\noindent 列女传：永平三年的东汉与边地、地方社会发生\eventanchor{EH-DENSE-0256}{以出入之勤，特封子成关内侯，官至齐相}。同卷前后条目所示的制度、地方或边疆条件构成其条件，作为跨材料核验的编年节点保留随之进入后续年序。《后汉书》·列女传提供这一线索，但《后汉书》传志兼具官修分类与后出叙事；地方实践、数字和群体名称须与简牍、碑刻及考古交叉核验。
\par

\noindent 在东夷列传：建武中元二年，\eventanchor{EH-DENSE-0257}{桓﹑灵闲，倭国大乱，更相攻伐，历年无主}把东汉与边地、地方社会置于同卷前后条目所示的制度、地方或边疆条件与作为跨材料核验的编年节点保留之间。可据《后汉书》·东夷列传追踪其大体经过；《后汉书》传志兼具官修分类与后出叙事；地方实践、数字和群体名称须与简牍、碑刻及考古交叉核验。
\par

\noindent 东汉与边地、地方社会在南蛮西南夷列传：和帝永元四年冬的\eventanchor{EH-DENSE-0258}{州郡募五里蛮六亭兵追击破之，皆散降}，并非孤立一事：同卷前后条目所示的制度、地方或边疆条件，且作为跨材料核验的编年节点保留。材料见《后汉书》·南蛮西南夷列传；《后汉书》传志兼具官修分类与后出叙事；地方实践、数字和群体名称须与简牍、碑刻及考古交叉核验。
\par

\noindent \eventanchor{EH-DENSE-0259}{林为下吏所欺，谬奏上滇岸以为大豪，承制封为归义侯，加号汉大都尉}发生于西羌传：永平元年，行动者为东汉与边地、地方社会。它回应同卷前后条目所示的制度、地方或边疆条件，并把作为跨材料核验的编年节点保留留给随后的人群和官署。《后汉书》·西羌传可作核验，惟《后汉书》传志兼具官修分类与后出叙事；地方实践、数字和群体名称须与简牍、碑刻及考古交叉核验。
\par

\noindent 从同卷前后条目所示的制度、地方或边疆条件出发，东汉与边地、地方社会于西域传：九年出现\eventanchor{EH-DENSE-0260}{北匈奴即复收属诸国，共为边寇十余岁}；作为跨材料核验的编年节点保留。这里采用《后汉书》·西域传的记载范围，不越过《后汉书》传志兼具官修分类与后出叙事；地方实践、数字和群体名称须与简牍、碑刻及考古交叉核验。
\par

\noindent 南匈奴列传：二十五年，东汉与边地、地方社会的\eventanchor{EH-DENSE-0261}{初，帝造战车，可驾数牛，上作楼橹，置于塞上，以拒匈奴}使同卷前后条目所示的制度、地方或边疆条件转为可见的行动，作为跨材料核验的编年节点保留。其依据为《后汉书》·南匈奴列传，《后汉书》传志兼具官修分类与后出叙事；地方实践、数字和群体名称须与简牍、碑刻及考古交叉核验。
\par

\noindent \eventanchor{EH-DENSE-0262}{匈奴国乱，乌桓乘弱击破之，匈奴转北徙数千里，漠南地空，帝乃以币帛赂乌桓}这一乌桓鲜卑列传：二十二年的节点，关联东汉与边地、地方社会、同卷前后条目所示的制度、地方或边疆条件和作为跨材料核验的编年节点保留。《后汉书》·乌桓鲜卑列传记录了它的基本轮廓；《后汉书》传志兼具官修分类与后出叙事；地方实践、数字和群体名称须与简牍、碑刻及考古交叉核验。
\par

\noindent 东汉与边地、地方社会在党锢列传：中平元年处理同卷前后条目所示的制度、地方或边疆条件时，出现\eventanchor{EH-DENSE-0263}{若久不赦宥，轻与张角合谋，为变滋大，悔之无救}；作为跨材料核验的编年节点保留。可核材料为《后汉书》·党锢列传，且《后汉书》传志兼具官修分类与后出叙事；地方实践、数字和群体名称须与简牍、碑刻及考古交叉核验。
\par

\noindent 循吏列传：建武二年的\eventanchor{EH-DENSE-0264}{流民稍还，渐成聚邑，使输租赋，同之平民}，让东汉与边地、地方社会的同卷前后条目所示的制度、地方或边疆条件在具体事务中显形；作为跨材料核验的编年节点保留。《后汉书》·循吏列传所能支持的范围仍受《后汉书》传志兼具官修分类与后出叙事；地方实践、数字和群体名称须与简牍、碑刻及考古交叉核验限制。
\par

\noindent 对东汉与边地、地方社会而言，\eventanchor{EH-DENSE-0265}{帝以经传之文多不正定，乃选通儒谒者刘珍及博士良史诣东观，各雠校*(汉)*家法，令伦监典其事}发生在宦者列传：四年：同卷前后条目所示的制度、地方或边疆条件，随后作为跨材料核验的编年节点保留。参看《后汉书》·宦者列传；《后汉书》传志兼具官修分类与后出叙事；地方实践、数字和群体名称须与简牍、碑刻及考古交叉核验。
\par

\noindent 同卷前后条目所示的制度、地方或边疆条件之下，东汉与边地、地方社会于儒林列传：建武五年经历\eventanchor{EH-DENSE-0266}{建初中，大会诸儒于白虎观，考详同异，连月乃罢}，并面对作为跨材料核验的编年节点保留。此处依《后汉书》·儒林列传，《后汉书》传志兼具官修分类与后出叙事；地方实践、数字和群体名称须与简牍、碑刻及考古交叉核验。
\par

\noindent 儒林列传：十二年，东汉与边地、地方社会面对同卷前后条目所示的制度、地方或边疆条件；\eventanchor{EH-DENSE-0267}{建武中，为州从事，征拜博士，稍迁太子少傅，卒于官}，其后作为跨材料核验的编年节点保留。记事依据《后汉书》·儒林列传；《后汉书》传志兼具官修分类与后出叙事；地方实践、数字和群体名称须与简牍、碑刻及考古交叉核验。
\par

\noindent 方术列传：八年的东汉与边地、地方社会发生\eventanchor{EH-DENSE-0268}{‘败我陂者翟子威，饴我大豆，亨我芋魁}。同卷前后条目所示的制度、地方或边疆条件构成其条件，作为跨材料核验的编年节点保留随之进入后续年序。《后汉书》·方术列传提供这一线索，但《后汉书》传志兼具官修分类与后出叙事；地方实践、数字和群体名称须与简牍、碑刻及考古交叉核验。
\par

\noindent 在方术列传：永建五年，\eventanchor{EH-DENSE-0269}{时暴风震雷，有声于外呼穆者三，穆不与语}把东汉与边地、地方社会置于同卷前后条目所示的制度、地方或边疆条件与作为跨材料核验的编年节点保留之间。可据《后汉书》·方术列传追踪其大体经过；《后汉书》传志兼具官修分类与后出叙事；地方实践、数字和群体名称须与简牍、碑刻及考古交叉核验。
\par

\noindent 东汉与边地、地方社会在列女传：永平三年的\eventanchor{EH-DENSE-0270}{时汉书始出，多未能通者，同郡马融伏于阁下，从昭受读，后又诏融兄续继昭成之}，并非孤立一事：同卷前后条目所示的制度、地方或边疆条件，且作为跨材料核验的编年节点保留。材料见《后汉书》·列女传；《后汉书》传志兼具官修分类与后出叙事；地方实践、数字和群体名称须与简牍、碑刻及考古交叉核验。
\par

\noindent \eventanchor{EH-DENSE-0271}{传言秦始皇遣方士徐福将童男女数千人入海，求蓬莱神仙不得，徐福畏诛不敢还，遂止此洲，世世相承，有数万家}发生于东夷列传：建武中元二年，行动者为东汉与边地、地方社会。它回应同卷前后条目所示的制度、地方或边疆条件，并把作为跨材料核验的编年节点保留留给随后的人群和官署。《后汉书》·东夷列传可作核验，惟《后汉书》传志兼具官修分类与后出叙事；地方实践、数字和群体名称须与简牍、碑刻及考古交叉核验。
\par

\noindent 从同卷前后条目所示的制度、地方或边疆条件出发，东汉与边地、地方社会于南蛮西南夷列传：和帝永元四年冬出现\eventanchor{EH-DENSE-0272}{顺帝永和元年，武陵太守上书，以蛮夷率服，可比汉人，增其租赋}；作为跨材料核验的编年节点保留。这里采用《后汉书》·南蛮西南夷列传的记载范围，不越过《后汉书》传志兼具官修分类与后出叙事；地方实践、数字和群体名称须与简牍、碑刻及考古交叉核验。
\par

\noindent 西羌传：永平元年，东汉与边地、地方社会的\eventanchor{EH-DENSE-0273}{明年，滇吾复降，林复奏其第一豪，与俱诣阙献见}使同卷前后条目所示的制度、地方或边疆条件转为可见的行动，作为跨材料核验的编年节点保留。其依据为《后汉书》·西羌传，《后汉书》传志兼具官修分类与后出叙事；地方实践、数字和群体名称须与简牍、碑刻及考古交叉核验。
\par

\noindent \eventanchor{EH-DENSE-0274}{敦煌太守曹宗患其暴害，元初六年，乃上遣行长史索班，将千余人屯伊吾以招抚之，于是车师前王及鄯善王来降}这一西域传：九年的节点，关联东汉与边地、地方社会、同卷前后条目所示的制度、地方或边疆条件和作为跨材料核验的编年节点保留。《后汉书》·西域传记录了它的基本轮廓；《后汉书》传志兼具官修分类与后出叙事；地方实践、数字和群体名称须与简牍、碑刻及考古交叉核验。
\par

\noindent 东汉与边地、地方社会在南匈奴列传：二十五年处理同卷前后条目所示的制度、地方或边疆条件时，出现\eventanchor{EH-DENSE-0275}{北部薁鞬骨都侯与右骨都侯率觽三万余人来归南单于，南单于复遣使诣阙，奉藩称臣，献国珍宝，求使者监护，遣侍子，修旧约}；作为跨材料核验的编年节点保留。可核材料为《后汉书》·南匈奴列传，且《后汉书》传志兼具官修分类与后出叙事；地方实践、数字和群体名称须与简牍、碑刻及考古交叉核验。
\par

\noindent 乌桓鲜卑列传：二十二年的\eventanchor{EH-DENSE-0276}{辽西乌桓大人郝旦等九百二十二人率觽向化，诣阙朝贡，献奴婢牛马及弓虎豹貂皮}，让东汉与边地、地方社会的同卷前后条目所示的制度、地方或边疆条件在具体事务中显形；作为跨材料核验的编年节点保留。《后汉书》·乌桓鲜卑列传所能支持的范围仍受《后汉书》传志兼具官修分类与后出叙事；地方实践、数字和群体名称须与简牍、碑刻及考古交叉核验限制。
\par

\noindent 对东汉与边地、地方社会而言，\eventanchor{EH-DENSE-0277}{”帝惧其言，乃大赦党人，诛徙之家皆归故郡}发生在党锢列传：中平元年：同卷前后条目所示的制度、地方或边疆条件，随后作为跨材料核验的编年节点保留。参看《后汉书》·党锢列传；《后汉书》传志兼具官修分类与后出叙事；地方实践、数字和群体名称须与简牍、碑刻及考古交叉核验。
\par

\noindent 同卷前后条目所示的制度、地方或边疆条件之下，东汉与边地、地方社会于循吏列传：建武二年经历\eventanchor{EH-DENSE-0278}{又耒阳县*(山)**[出]*铁石，佗郡民庶常依因聚会，私为冶铸，遂招来亡命，多致奸盗}，并面对作为跨材料核验的编年节点保留。此处依《后汉书》·循吏列传，《后汉书》传志兼具官修分类与后出叙事；地方实践、数字和群体名称须与简牍、碑刻及考古交叉核验。
\par

\noindent 宦者列传：四年，东汉与边地、地方社会面对同卷前后条目所示的制度、地方或边疆条件；\eventanchor{EH-DENSE-0279}{小黄门李闰与帝乳母王圣常共谮太后兄执金吾悝等，言欲废帝，立平原王*(德)**[翼]*，帝每忿惧}，其后作为跨材料核验的编年节点保留。记事依据《后汉书》·宦者列传；《后汉书》传志兼具官修分类与后出叙事；地方实践、数字和群体名称须与简牍、碑刻及考古交叉核验。
\par

\noindent 儒林列传：建武五年的东汉与边地、地方社会发生\eventanchor{EH-DENSE-0280}{又诏高才生受古文尚书、毛诗、谷梁、左氏春秋，虽不立学官，然皆擢高第为讲郎，给事近署，所以网罗遗逸，博存觽家}。同卷前后条目所示的制度、地方或边疆条件构成其条件，作为跨材料核验的编年节点保留随之进入后续年序。《后汉书》·儒林列传提供这一线索，但《后汉书》传志兼具官修分类与后出叙事；地方实践、数字和群体名称须与简牍、碑刻及考古交叉核验。
\par

\noindent 在儒林列传：十二年，\eventanchor{EH-DENSE-0281}{建武中，赵节王栩闻其高名，遣使赍玉帛请以为师，望不受}把东汉与边地、地方社会置于同卷前后条目所示的制度、地方或边疆条件与作为跨材料核验的编年节点保留之间。可据《后汉书》·儒林列传追踪其大体经过；《后汉书》传志兼具官修分类与后出叙事；地方实践、数字和群体名称须与简牍、碑刻及考古交叉核验。
\par

\noindent 东汉与边地、地方社会在方术列传：八年的\eventanchor{EH-DENSE-0282}{’昔大禹决江疏河以利天下，明府今兴立废业，富国安民，童谣之言，将有征于此}，并非孤立一事：同卷前后条目所示的制度、地方或边疆条件，且作为跨材料核验的编年节点保留。材料见《后汉书》·方术列传；《后汉书》传志兼具官修分类与后出叙事；地方实践、数字和群体名称须与简牍、碑刻及考古交叉核验。
\par

\noindent \eventanchor{EH-DENSE-0283}{有顷，呼者自牖而入，音状甚怪，穆诵经自若，终亦无它妖异，时人奇之}发生于方术列传：永建五年，行动者为东汉与边地、地方社会。它回应同卷前后条目所示的制度、地方或边疆条件，并把作为跨材料核验的编年节点保留留给随后的人群和官署。《后汉书》·方术列传可作核验，惟《后汉书》传志兼具官修分类与后出叙事；地方实践、数字和群体名称须与简牍、碑刻及考古交叉核验。
\par

\noindent 从同卷前后条目所示的制度、地方或边疆条件出发，东汉与边地、地方社会于列女传：永平三年出现\eventanchor{EH-DENSE-0284}{妾闻谦让之风，德莫大焉，故典坟述美，神只降福}；作为跨材料核验的编年节点保留。这里采用《后汉书》·列女传的记载范围，不越过《后汉书》传志兼具官修分类与后出叙事；地方实践、数字和群体名称须与简牍、碑刻及考古交叉核验。
\par

\noindent 东夷列传：建武中元二年，东汉与边地、地方社会的\eventanchor{EH-DENSE-0285}{会稽东冶县人有入海行遭风，流移至澶洲者}使同卷前后条目所示的制度、地方或边疆条件转为可见的行动，作为跨材料核验的编年节点保留。其依据为《后汉书》·东夷列传，《后汉书》传志兼具官修分类与后出叙事；地方实践、数字和群体名称须与简牍、碑刻及考古交叉核验。
\par

\noindent \eventanchor{EH-DENSE-0286}{是故羇縻而绥抚之，附则受而不逆，叛则□而不追}这一南蛮西南夷列传：和帝永元四年冬的节点，关联东汉与边地、地方社会、同卷前后条目所示的制度、地方或边疆条件和作为跨材料核验的编年节点保留。《后汉书》·南蛮西南夷列传记录了它的基本轮廓；《后汉书》传志兼具官修分类与后出叙事；地方实践、数字和群体名称须与简牍、碑刻及考古交叉核验。
\par

\noindent 东汉与边地、地方社会在西羌传：永平元年处理同卷前后条目所示的制度、地方或边疆条件时，出现\eventanchor{EH-DENSE-0287}{滇吾子东吾立，以父降汉，乃入居塞内，谨愿自守}；作为跨材料核验的编年节点保留。可核材料为《后汉书》·西羌传，且《后汉书》传志兼具官修分类与后出叙事；地方实践、数字和群体名称须与简牍、碑刻及考古交叉核验。
\par

\noindent 西域传：九年的\eventanchor{EH-DENSE-0288}{数月，北匈奴复率车师后部王共攻没班等，遂击走其前王}，让东汉与边地、地方社会的同卷前后条目所示的制度、地方或边疆条件在具体事务中显形；作为跨材料核验的编年节点保留。《后汉书》·西域传所能支持的范围仍受《后汉书》传志兼具官修分类与后出叙事；地方实践、数字和群体名称须与简牍、碑刻及考古交叉核验限制。
\par

\noindent 对东汉与边地、地方社会而言，\eventanchor{EH-DENSE-0289}{遣中郎将段郴、副校尉王郁使南单于，立其庭，去五原西部塞八十里}发生在南匈奴列传：二十六年：同卷前后条目所示的制度、地方或边疆条件，随后作为跨材料核验的编年节点保留。参看《后汉书》·南匈奴列传；《后汉书》传志兼具官修分类与后出叙事；地方实践、数字和群体名称须与简牍、碑刻及考古交叉核验。
\par

\noindent 同卷前后条目所示的制度、地方或边疆条件之下，东汉与边地、地方社会于乌桓鲜卑列传：二十二年经历\eventanchor{EH-DENSE-0290}{是时四夷朝贺，络驿而至，天子乃命大会劳飨，赐以珍宝}，并面对作为跨材料核验的编年节点保留。此处依《后汉书》·乌桓鲜卑列传，《后汉书》传志兼具官修分类与后出叙事；地方实践、数字和群体名称须与简牍、碑刻及考古交叉核验。
\par

\noindent 党锢列传：中平元年，东汉与边地、地方社会面对同卷前后条目所示的制度、地方或边疆条件；\eventanchor{EH-DENSE-0291}{其后黄巾遂盛，朝野崩离，纲纪文章荡然矣}，其后作为跨材料核验的编年节点保留。记事依据《后汉书》·党锢列传；《后汉书》传志兼具官修分类与后出叙事；地方实践、数字和群体名称须与简牍、碑刻及考古交叉核验。
\par

\noindent 循吏列传：建武二年的东汉与边地、地方社会发生\eventanchor{EH-DENSE-0292}{飒乃上起铁官，罢斥私铸，岁所增入五百余万}。同卷前后条目所示的制度、地方或边疆条件构成其条件，作为跨材料核验的编年节点保留随之进入后续年序。《后汉书》·循吏列传提供这一线索，但《后汉书》传志兼具官修分类与后出叙事；地方实践、数字和群体名称须与简牍、碑刻及考古交叉核验。
\par

\noindent 在宦者列传：四年，\eventanchor{EH-DENSE-0293}{及太后崩，遂诛邓氏而废平原王，封闰雍乡侯；又小黄门江京以谗谄进，初迎帝于邸，以功封都乡侯，食邑各三百户}把东汉与边地、地方社会置于同卷前后条目所示的制度、地方或边疆条件与作为跨材料核验的编年节点保留之间。可据《后汉书》·宦者列传追踪其大体经过；《后汉书》传志兼具官修分类与后出叙事；地方实践、数字和群体名称须与简牍、碑刻及考古交叉核验。
\par

\noindent 东汉与边地、地方社会在儒林列传：建武五年的\eventanchor{EH-DENSE-0294}{时樊准、徐防并陈敦学之宜，又言儒职多非其人，于是制诏公卿妙简其选，三署郎能通经术者，皆得察举}，并非孤立一事：同卷前后条目所示的制度、地方或边疆条件，且作为跨材料核验的编年节点保留。材料见《后汉书》·儒林列传；《后汉书》传志兼具官修分类与后出叙事；地方实践、数字和群体名称须与简牍、碑刻及考古交叉核验。
\par

\noindent \eventanchor{EH-DENSE-0295}{永平初，为侍中﹑越骑校尉，入讲省内}发生于儒林列传：十二年，行动者为东汉与边地、地方社会。它回应同卷前后条目所示的制度、地方或边疆条件，并把作为跨材料核验的编年节点保留留给随后的人群和官署。《后汉书》·儒林列传可作核验，惟《后汉书》传志兼具官修分类与后出叙事；地方实践、数字和群体名称须与简牍、碑刻及考古交叉核验。
\par

\noindent 从同卷前后条目所示的制度、地方或边疆条件出发，东汉与边地、地方社会于方术列传：八年出现\eventanchor{EH-DENSE-0296}{”晨大悦，因署杨为都水掾，使典其事}；作为跨材料核验的编年节点保留。这里采用《后汉书》·方术列传的记载范围，不越过《后汉书》传志兼具官修分类与后出叙事；地方实践、数字和群体名称须与简牍、碑刻及考古交叉核验。
\par

\noindent 方术列传：永建五年，东汉与边地、地方社会的\eventanchor{EH-DENSE-0297}{夫富贵在天，得之有命，以货求位，吾不忍也}使同卷前后条目所示的制度、地方或边疆条件转为可见的行动，作为跨材料核验的编年节点保留。其依据为《后汉书》·方术列传，《后汉书》传志兼具官修分类与后出叙事；地方实践、数字和群体名称须与简牍、碑刻及考古交叉核验。
\par

\noindent \eventanchor{EH-DENSE-0298}{夙夜劬心，勤不告劳，而今而后，乃知免耳}这一列女传：永平三年的节点，关联东汉与边地、地方社会、同卷前后条目所示的制度、地方或边疆条件和作为跨材料核验的编年节点保留。《后汉书》·列女传记录了它的基本轮廓；《后汉书》传志兼具官修分类与后出叙事；地方实践、数字和群体名称须与简牍、碑刻及考古交叉核验。
\par

\noindent 东汉与边地、地方社会在南蛮西南夷列传：和帝永元四年冬处理同卷前后条目所示的制度、地方或边疆条件时，出现\eventanchor{EH-DENSE-0299}{先帝旧典，贡税多少，所由来久矣}；作为跨材料核验的编年节点保留。可核材料为《后汉书》·南蛮西南夷列传，且《后汉书》传志兼具官修分类与后出叙事；地方实践、数字和群体名称须与简牍、碑刻及考古交叉核验。
\par

\noindent 西羌传：永平元年的\eventanchor{EH-DENSE-0300}{肃宗建初元年，安夷县吏略妻卑湳种羌妇，吏为其夫所杀，安夷长宗延追之出塞，种人恐见诛，遂共杀延，而与勒姐及吾良二种相结为寇}，让东汉与边地、地方社会的同卷前后条目所示的制度、地方或边疆条件在具体事务中显形；作为跨材料核验的编年节点保留。《后汉书》·西羌传所能支持的范围仍受《后汉书》传志兼具官修分类与后出叙事；地方实践、数字和群体名称须与简牍、碑刻及考古交叉核验限制。
\par

\noindent 对东汉与边地、地方社会而言，\eventanchor{EH-DENSE-0301}{鄯善逼急，求救于曹宗，宗因此请出兵击匈奴，报索班之耻，复欲进取西域}发生在西域传：九年：同卷前后条目所示的制度、地方或边疆条件，随后作为跨材料核验的编年节点保留。参看《后汉书》·西域传；《后汉书》传志兼具官修分类与后出叙事；地方实践、数字和群体名称须与简牍、碑刻及考古交叉核验。
\par

\noindent 同卷前后条目所示的制度、地方或边疆条件之下，东汉与边地、地方社会于南匈奴列传：二十六年经历\eventanchor{EH-DENSE-0302}{郴等反命，诏乃听南单于入居云中}，并面对作为跨材料核验的编年节点保留。此处依《后汉书》·南匈奴列传，《后汉书》传志兼具官修分类与后出叙事；地方实践、数字和群体名称须与简牍、碑刻及考古交叉核验。
\par

\noindent 乌桓鲜卑列传：二十二年，东汉与边地、地方社会面对同卷前后条目所示的制度、地方或边疆条件；\eventanchor{EH-DENSE-0303}{乌桓或愿留宿韂，于是封其渠帅为侯王君长者八十一人，皆居塞内，布于缘边诸郡，令招来种人，给其衣食，遂为汉侦候，助击匈奴、鲜卑}，其后作为跨材料核验的编年节点保留。记事依据《后汉书》·乌桓鲜卑列传；《后汉书》传志兼具官修分类与后出叙事；地方实践、数字和群体名称须与简牍、碑刻及考古交叉核验。
\par

\noindent 党锢列传：中平元年的东汉与边地、地方社会发生\eventanchor{EH-DENSE-0304}{翟超，山阳太守，事见陈蕃传，字及郡县未详}。同卷前后条目所示的制度、地方或边疆条件构成其条件，作为跨材料核验的编年节点保留随之进入后续年序。《后汉书》·党锢列传提供这一线索，但《后汉书》传志兼具官修分类与后出叙事；地方实践、数字和群体名称须与简牍、碑刻及考古交叉核验。
\par

\noindent 在循吏列传：二十五年，\eventanchor{EH-DENSE-0305}{光武欲以为少府，会飒被疾，不能拜起，□以桂阳太守归家，须后诏书}把东汉与边地、地方社会置于同卷前后条目所示的制度、地方或边疆条件与作为跨材料核验的编年节点保留之间。可据《后汉书》·循吏列传追踪其大体经过；《后汉书》传志兼具官修分类与后出叙事；地方实践、数字和群体名称须与简牍、碑刻及考古交叉核验。
\par

\noindent 东汉与边地、地方社会在宦者列传：四年的\eventanchor{EH-DENSE-0306}{闰﹑京并迁中常侍，江京兼大长秋，与中常侍樊丰﹑黄门令刘安﹑钩盾令陈达及王圣﹑圣女伯荣扇动内外，竞为侈虐}，并非孤立一事：同卷前后条目所示的制度、地方或边疆条件，且作为跨材料核验的编年节点保留。材料见《后汉书》·宦者列传；《后汉书》传志兼具官修分类与后出叙事；地方实践、数字和群体名称须与简牍、碑刻及考古交叉核验。
\par

\noindent \eventanchor{EH-DENSE-0307}{顺帝感翟酺之言，乃更修黉宇，凡所造构二百四十房，千八百五十室}发生于儒林列传：建武五年，行动者为东汉与边地、地方社会。它回应同卷前后条目所示的制度、地方或边疆条件，并把作为跨材料核验的编年节点保留留给随后的人群和官署。《后汉书》·儒林列传可作核验，惟《后汉书》传志兼具官修分类与后出叙事；地方实践、数字和群体名称须与简牍、碑刻及考古交叉核验。
\par

\noindent 从同卷前后条目所示的制度、地方或边疆条件出发，东汉与边地、地方社会于儒林列传：建初五年出现\eventanchor{EH-DENSE-0308}{建初三年，举孝廉，迁海西令，卒于官}；作为跨材料核验的编年节点保留。这里采用《后汉书》·儒林列传的记载范围，不越过《后汉书》传志兼具官修分类与后出叙事；地方实践、数字和群体名称须与简牍、碑刻及考古交叉核验。
\par

\noindent 方术列传：八年，东汉与边地、地方社会的\eventanchor{EH-DENSE-0309}{杨因高下形埶，起塘四百余里，数年乃立}使同卷前后条目所示的制度、地方或边疆条件转为可见的行动，作为跨材料核验的编年节点保留。其依据为《后汉书》·方术列传，《后汉书》传志兼具官修分类与后出叙事；地方实践、数字和群体名称须与简牍、碑刻及考古交叉核验。
\par

\noindent \eventanchor{EH-DENSE-0310}{后举孝廉，以高第为主事，迁缯相}这一方术列传：永建五年的节点，关联东汉与边地、地方社会、同卷前后条目所示的制度、地方或边疆条件和作为跨材料核验的编年节点保留。《后汉书》·方术列传记录了它的基本轮廓；《后汉书》传志兼具官修分类与后出叙事；地方实践、数字和群体名称须与简牍、碑刻及考古交叉核验。
\par

\noindent 东汉与边地、地方社会在列女传：永平三年处理同卷前后条目所示的制度、地方或边疆条件时，出现\eventanchor{EH-DENSE-0311}{吾性疏顽，教道无素，恒恐子谷负辱清朝}；作为跨材料核验的编年节点保留。可核材料为《后汉书》·列女传，且《后汉书》传志兼具官修分类与后出叙事；地方实践、数字和群体名称须与简牍、碑刻及考古交叉核验。
\par

\noindent 南蛮西南夷列传：和帝永元四年冬的\eventanchor{EH-DENSE-0312}{其冬澧中、溇中蛮果争贡布非旧约，遂杀乡吏，举种反叛}，让东汉与边地、地方社会的同卷前后条目所示的制度、地方或边疆条件在具体事务中显形；作为跨材料核验的编年节点保留。《后汉书》·南蛮西南夷列传所能支持的范围仍受《后汉书》传志兼具官修分类与后出叙事；地方实践、数字和群体名称须与简牍、碑刻及考古交叉核验限制。
\par

\noindent 对东汉与边地、地方社会而言，\eventanchor{EH-DENSE-0313}{陇西太守孙纯遣从事李睦及金城兵会和罗谷，与卑湳等战，斩首虏数百人}发生在西羌传：永平元年：同卷前后条目所示的制度、地方或边疆条件，随后作为跨材料核验的编年节点保留。参看《后汉书》·西羌传；《后汉书》传志兼具官修分类与后出叙事；地方实践、数字和群体名称须与简牍、碑刻及考古交叉核验。
\par

\noindent 同卷前后条目所示的制度、地方或边疆条件之下，东汉与边地、地方社会于西域传：九年经历\eventanchor{EH-DENSE-0314}{邓太后不许，但令置护西域副校尉，居敦煌，复部营兵三百人，羁縻而已}，并面对作为跨材料核验的编年节点保留。此处依《后汉书》·西域传，《后汉书》传志兼具官修分类与后出叙事；地方实践、数字和群体名称须与简牍、碑刻及考古交叉核验。
\par

\noindent 南匈奴列传：二十六年，东汉与边地、地方社会面对同卷前后条目所示的制度、地方或边疆条件；\eventanchor{EH-DENSE-0315}{遣使上书，献骆喰二头，文马十匹}，其后作为跨材料核验的编年节点保留。记事依据《后汉书》·南匈奴列传；《后汉书》传志兼具官修分类与后出叙事；地方实践、数字和群体名称须与简牍、碑刻及考古交叉核验。
\par

\noindent 乌桓鲜卑列传：二十二年的东汉与边地、地方社会发生\eventanchor{EH-DENSE-0316}{时司徒掾班彪上言：“乌桓天性轻黠，好为寇贼，若久放纵而无总领者，必复侵掠居人，但委主降掾史，恐非所能制}。同卷前后条目所示的制度、地方或边疆条件构成其条件，作为跨材料核验的编年节点保留随之进入后续年序。《后汉书》·乌桓鲜卑列传提供这一线索，但《后汉书》传志兼具官修分类与后出叙事；地方实践、数字和群体名称须与简牍、碑刻及考古交叉核验。
\par

\noindent 在党锢列传：中平元年，\eventanchor{EH-DENSE-0317}{淑少学明五经，遂隐居，立精舍讲授，诸生常数百人}把东汉与边地、地方社会置于同卷前后条目所示的制度、地方或边疆条件与作为跨材料核验的编年节点保留之间。可据《后汉书》·党锢列传追踪其大体经过；《后汉书》传志兼具官修分类与后出叙事；地方实践、数字和群体名称须与简牍、碑刻及考古交叉核验。
\par

\noindent 东汉与边地、地方社会在循吏列传：二十五年的\eventanchor{EH-DENSE-0318}{居二岁，载病诣阙，自陈困笃，乃收印绶，赐钱十万，后卒于家}，并非孤立一事：同卷前后条目所示的制度、地方或边疆条件，且作为跨材料核验的编年节点保留。材料见《后汉书》·循吏列传；《后汉书》传志兼具官修分类与后出叙事；地方实践、数字和群体名称须与简牍、碑刻及考古交叉核验。
\par

\noindent \eventanchor{EH-DENSE-0319}{又帝舅大将军耿宝﹑皇后兄大鸿胪阎显更相阿党，遂枉杀太尉杨震，废皇太子为济阴王}发生于宦者列传：四年，行动者为东汉与边地、地方社会。它回应同卷前后条目所示的制度、地方或边疆条件，并把作为跨材料核验的编年节点保留留给随后的人群和官署。《后汉书》·宦者列传可作核验，惟《后汉书》传志兼具官修分类与后出叙事；地方实践、数字和群体名称须与简牍、碑刻及考古交叉核验。
\par

\noindent 从同卷前后条目所示的制度、地方或边疆条件出发，东汉与边地、地方社会于儒林列传：建武五年出现\eventanchor{EH-DENSE-0320}{试明经下第补弟子，增甲乙之科员各十人，除郡国耆儒皆补郎、舍人}；作为跨材料核验的编年节点保留。这里采用《后汉书》·儒林列传的记载范围，不越过《后汉书》传志兼具官修分类与后出叙事；地方实践、数字和群体名称须与简牍、碑刻及考古交叉核验。
\par

\noindent 儒林列传：建初五年，东汉与边地、地方社会的\eventanchor{EH-DENSE-0321}{建武初，举明经，补弘农文学，迁陈仓县丞}使同卷前后条目所示的制度、地方或边疆条件转为可见的行动，作为跨材料核验的编年节点保留。其依据为《后汉书》·儒林列传，《后汉书》传志兼具官修分类与后出叙事；地方实践、数字和群体名称须与简牍、碑刻及考古交叉核验。
\par

\noindent \eventanchor{EH-DENSE-0322}{太守闻忠信可以感灵，今其效乎！”即夜出杨，遣归}这一方术列传：八年的节点，关联东汉与边地、地方社会、同卷前后条目所示的制度、地方或边疆条件和作为跨材料核验的编年节点保留。《后汉书》·方术列传记录了它的基本轮廓；《后汉书》传志兼具官修分类与后出叙事；地方实践、数字和群体名称须与简牍、碑刻及考古交叉核验。
\par

\noindent 东汉与边地、地方社会在方术列传：永建五年处理同卷前后条目所示的制度、地方或边疆条件时，出现\eventanchor{EH-DENSE-0323}{时缯侯刘敞，东海恭王之后也，所为多不法，废嫡立庶，傲很放恣}；作为跨材料核验的编年节点保留。可核材料为《后汉书》·方术列传，且《后汉书》传志兼具官修分类与后出叙事；地方实践、数字和群体名称须与简牍、碑刻及考古交叉核验。
\par

\noindent 列女传：永平三年的\eventanchor{EH-DENSE-0324}{正色端操，以事夫主，清静自守，无好戏笑，絜齐酒食，以供祖宗，是谓继祭祀也}，让东汉与边地、地方社会的同卷前后条目所示的制度、地方或边疆条件在具体事务中显形；作为跨材料核验的编年节点保留。《后汉书》·列女传所能支持的范围仍受《后汉书》传志兼具官修分类与后出叙事；地方实践、数字和群体名称须与简牍、碑刻及考古交叉核验限制。
\par

\noindent 对东汉与边地、地方社会而言，\eventanchor{EH-DENSE-0325}{明年春，蛮二万人围充城，八千人寇夷道}发生在南蛮西南夷列传：和帝永元四年冬：同卷前后条目所示的制度、地方或边疆条件，随后作为跨材料核验的编年节点保留。参看《后汉书》·南蛮西南夷列传；《后汉书》传志兼具官修分类与后出叙事；地方实践、数字和群体名称须与简牍、碑刻及考古交叉核验。
\par

\noindent 同卷前后条目所示的制度、地方或边疆条件之下，东汉与边地、地方社会于西羌传：永平元年经历\eventanchor{EH-DENSE-0326}{复拜故度辽将军吴棠领护羌校尉，居安夷}，并面对作为跨材料核验的编年节点保留。此处依《后汉书》·西羌传，《后汉书》传志兼具官修分类与后出叙事；地方实践、数字和群体名称须与简牍、碑刻及考古交叉核验。
\par

\noindent 西域传：九年，东汉与边地、地方社会面对同卷前后条目所示的制度、地方或边疆条件；\eventanchor{EH-DENSE-0327}{其后北虏连与车师入寇河西，朝廷不能禁，议者因欲闭玉门、阳关，以绝其患}，其后作为跨材料核验的编年节点保留。记事依据《后汉书》·西域传；《后汉书》传志兼具官修分类与后出叙事；地方实践、数字和群体名称须与简牍、碑刻及考古交叉核验。
\par

\noindent 南匈奴列传：二十六年的东汉与边地、地方社会发生\eventanchor{EH-DENSE-0328}{夏，南单于所获北虏薁鞬左贤王将其觽及南部五骨都侯合三万余人畔归，去北庭三百余里，共立薁鞬左贤王为单于}。同卷前后条目所示的制度、地方或边疆条件构成其条件，作为跨材料核验的编年节点保留随之进入后续年序。《后汉书》·南匈奴列传提供这一线索，但《后汉书》传志兼具官修分类与后出叙事；地方实践、数字和群体名称须与简牍、碑刻及考古交叉核验。
\par
